\documentclass[aspectratio=169,11pt]{beamer}
% For Chinese slides: \documentclass[aspectratio=169,11pt]{ctexbeamer}, compiled with XeLaTeX.

\usepackage{amsmath,amssymb}
% Add \usepackage{graphicx} or pgfplots only when the deck includes figures; call figures by relative path, e.g. \includegraphics[width=\linewidth]{figures/fig1.png} (no \graphicspath).

\setbeamertemplate{navigation symbols}{}
\setbeamertemplate{footline}[frame number]

\title{Paper Title}
\subtitle{Author (Year), Journal}
\author{Presenter}
\date{}

\begin{document}

\begin{frame}
  \titlepage
\end{frame}

\begin{frame}{Research Question}
  \begin{itemize}
    \item The question the paper asks
    \item Why it matters
    \item The answer, previewed
  \end{itemize}
\end{frame}

\begin{frame}{Core Mechanism}
  \[
    \alert{Y} = F(X_1, X_2)
  \]
  \begin{itemize}
    \item Define each symbol where it first appears
    \item One sentence of intuition
  \end{itemize}
\end{frame}

\begin{frame}{Takeaways}
  \begin{itemize}
    \item One line per takeaway
  \end{itemize}
\end{frame}

\begin{frame}{Discussion}
  \begin{itemize}
    \item Which assumption drives the result?
  \end{itemize}
\end{frame}

\appendix

\begin{frame}{Backup: derivation}
  Multi-step derivations go here.
\end{frame}

\end{document}
